\paragraph{Nguyên tắc giao thoa nhiều thành phần đơn sắc}
\begin{itemize}
\item Hiện tượng giao thoa ánh sáng chỉ xảy ra nếu ánh sáng có cùng bước sóng và độ lệch pha không đổi (thông thường người ta tạo ra hai nguồn thứ cấp $S_1$ và $S_2$ cùng pha). 
\item Một nguồn sáng phức tạp gồm nhiều thành phần đơn sắc (bước sóng khác nhau), các thành phần đơn sắc khác nhau không thể giao thoa với nhau. 
\item Khi bài toán đề cập đến việc giao thoa với nhiều bức xạ hay giao thoa của ánh sáng trắng đó là hiện tượng giao thoa của từng bức xạ đơn sắc \bth{\Rightarrow} tạo nên hệ các vân sáng, vân tối riêng biệt chồng chất lên nhau.
\end{itemize}
\paragraph{Giao thoa với hai bức xạ đơn sắc}
\Binhluan{
\begin{itemize}
\item Khi nguồn sáng S phát ra hai bức xạ (ánh sáng) có bước sóng \bth{\lambda_1} và \bth{\lambda_2} thì:
\begin{itemize}
\item Trên màn có hai hệ vân giao thoa ứng với hai ánh sáng có bước sóng \bth{\lambda_1} và \bth{\lambda_2}.
\item Ở vị trí trung tâm có hai vân sáng trung tâm trùng nhau
\item Tại các vị trí M, N,\ldots thì hai vân lại trùng nhau khi 
\begin{align*}
&x_{s1}=x_{s2}\Leftrightarrow k_1i_1=k_2i_2 \Leftrightarrow k_1\lambda_1=k_2\lambda_2\\ 
\Rightarrow &\boder{\dfrac{i_2}{i_1}=\dfrac{\lambda_2}{\lambda_1}=\dfrac{b}{c}}\ \left(\ct{phân số tối giản}\right)
\end{align*}
\end{itemize}
\item Khoảng vân trùng (khoảng cách nhỏ nhất giữa hai vân cùng màu với vân trung tâm)
\begin{align*}
\boder{i_{\equiv}=bi_1=ci_2}\ \ct{hoặc}\ i_{12}=BCNN(i_1,i_2)
\end{align*}
\end{itemize}
}
{Tại O luôn luôn là một vị trí trùng nhau.\\
Về mặt định lượng ta có thể coi như hệ vân giao thoa với một bức xạ "mới" có bước sóng $\lambda_{\equiv}$ và khoảng vân $i_{\equiv}$.
}